\documentclass[11pt]{article}
\usepackage[margin=1in]{geometry}
\usepackage{amsmath, amssymb, amsthm, xcolor, mathpartir}
\usepackage{hyperref}

\newtheorem{task}{Task}
\newtheorem{lemma}{Lemma}
\newtheorem{thm}{Theorem}
  

%% Inference rules
%Inference rule with a title
\newcommand{\Rule}[4][]{\ensuremath{\inferrule*[right={(#2)},#1]{#3}{#4}}}
%Inference rule with no title
\newcommand{\RuleNoLabel}[3][]{\ensuremath{\inferrule*[#1]{#2}{#3}}}

%% BNF
\newcommand{\bnfdef}{\mathop{::=}}
\newcommand{\bnfalt}{\mathbin{|}}

%% XML
\newcommand{\isxml}[1]{#1~\ensuremath{\mathsf{is~XML}}}
\newcommand{\xml}[1]{\text{\texttt{#1}}}

%% RB Trees
\newcommand{\tree}{\ensuremath{T}}
\newcommand{\leaf}{\ensuremath{\mathsf{Leaf}}}
\newcommand{\bnode}[2]{\ensuremath{\mathsf{BNode}(#1, #2)}}
\newcommand{\rnode}[2]{\ensuremath{{\color{red}{\mathsf{RNode}}}(#1, #2)}}
\newcommand{\bc}{\ensuremath{\mathsf{Black}}}
\newcommand{\rc}{\ensuremath{{\color{red}{\mathsf{Red}}}}}
\newcommand{\istree}[2]{~\ensuremath{\mathsf{isTree}~#1~#2}}
\newcommand{\colr}{\ensuremath{c}}
\newcommand{\nodes}[1]{\ensuremath{\mathit{Nodes}(#1)}}


% Shorthand for the WORKED EXAMPLES (not the homework's constructs)
\newcommand{\val}[1]{#1\ \mathsf{val}}
\newcommand{\stepto}[2]{#1 \mapsto #2}
\newcommand{\PAIR}[2]{(#1,\, #2)}
\newcommand{\FST}[1]{\mathsf{fst}\ #1}
\newcommand{\SQRT}[1]{\mathsf{sqrt}\ #1}
\newcommand{\INT}{\mathsf{int}}
\newcommand{\STRING}{\mathsf{string}}
\newcommand{\NN}{\mathsf{nn}}
\newcommand{\ANY}{\mathsf{any}}
\newcommand{\NP}{\mathsf{np}}
\newcommand{\NZ}{\mathsf{nz}}
\newcommand{\MZ}{\mathsf{mz}}
\newcommand{\num}[1]{\overline{#1}}
\newcommand{\str}[1]{\text{``#1''}}
\newcommand{\blank}{\underline{\hspace{1.5cm}}}

\title{CSE5095 --- Homework 2 Walkthrough\\
\large Type Safety, by Worked Analogy}
\author{}
\date{}

\begin{document}
\maketitle

\textbf{How to use this.} Every section below works a \emph{different}
construct than the homework, but uses \emph{the same technique} you need.
Read the worked example, then do the matching task yourself (the
fill-in-the-blank skeletons in \texttt{hw2\_hints.tex} pair well with this).
Nothing here is an answer to a homework task. Your syllabus allows AI for
understanding the material; the submitted write-up has to be yours.

\paragraph{The four judgments you'll use.}
\begin{itemize}
  \item $\val{e}$: $e$ is finished evaluating.
  \item $\stepto{e}{e'}$: $e$ takes one step to $e'$.
  \item $e:\tau$: $e$ has type $\tau$.
  \item Theorems: \textbf{Progress} (if $e:\tau$ then $\val{e}$ or
        $\exists e'.\ \stepto{e}{e'}$) and \textbf{Preservation} (if $e:\tau$
        and $\stepto{e}{e'}$ then $e':\tau$).
\end{itemize}

\paragraph{The one distinction to keep in your head.}
\begin{itemize}
  \item A \emph{search rule} has a \emph{step} judgment as a premise
        (e.g.\ StepSearchAddLeft: ``if $e_1\mapsto e_1'$ then
        $e_1+e_2\mapsto e_1'+e_2$''). It just says \emph{where} to evaluate
        next.
  \item A \emph{primitive / reduction rule} (StepAdd, StepCat,\dots) does
        the actual computation. It has no step premises, though it may have
        side conditions like $n_2\neq 0$.
\end{itemize}
In proofs, reduction-rule cases use \textbf{inversion}. Search-rule cases use
\textbf{the induction hypothesis}.

% ====================================================================
\part*{Part A: Booleans (Tasks 1--6)}
\addcontentsline{toc}{part}{Part A}

\textbf{Worked construct: pairs.} We extend the class language $\mathsf{E}$ with
\[
  \tau \bnfdef \cdots \bnfalt \tau_1\times\tau_2
  \qquad\qquad
  e \bnfdef \cdots \bnfalt \PAIR{e_1}{e_2} \bnfalt \FST{e}
\]
Intended meaning: evaluate both components left to right; $\FST{}$ of a
finished pair returns the first component.

% --------------------------------------------------------------------
\section*{A1. Writing dynamic rules \hfill (technique for Task 1)}

\paragraph{Recipe.}
\begin{enumerate}
  \item \emph{Value rules:} which finished forms of the new syntax are done?
  \item \emph{Reduction rules:} once the sub-expressions you care about are
        values, what does the construct compute to?
  \item \emph{Search rules:} for each sub-expression that must be evaluated
        \emph{before} reducing, add one rule that steps it and leaves
        everything else alone. Do \textbf{not} add search rules for
        sub-expressions that shouldn't be evaluated yet.
\end{enumerate}

\paragraph{Worked.} A pair is a value when both parts are values. This value
rule \emph{has premises}, unlike ValNum:
\[
  \Rule{ValPair}{\val{e_1} \\ \val{e_2}}{\val{\PAIR{e_1}{e_2}}}
\]
Reduction: $\FST{}$ of a finished pair gives the first component:
\[
  \Rule{StepFst}{\val{e_1} \\ \val{e_2}}{\stepto{\FST{\PAIR{e_1}{e_2}}}{e_1}}
\]
Search: we need to evaluate the left component, then the right one (only
once the left is a value), and the argument of $\FST{}$:
\[
  \Rule{StepSearchPairL}{\stepto{e_1}{e_1'}}
       {\stepto{\PAIR{e_1}{e_2}}{\PAIR{e_1'}{e_2}}}
  \quad
  \Rule{StepSearchPairR}{\val{e_1} \\ \stepto{e_2}{e_2'}}
       {\stepto{\PAIR{e_1}{e_2}}{\PAIR{e_1}{e_2'}}}
\]
\[
  \Rule{StepSearchFst}{\stepto{e}{e'}}{\stepto{\FST{e}}{\FST{e'}}}
\]

\paragraph{Sanity checks to run on your own rules.}
\begin{itemize}
  \item For every non-value expression, is there \emph{exactly one} rule
        that applies? (No rule gives a stuck expression; two rules
        means your semantics isn't deterministic.)
  \item Does anything get evaluated that the English description says
        \emph{shouldn't} be? For $\mathsf{if}$, reread the assignment's
        $\mathsf{if}\ x = 0\ \mathsf{then}\ 0\ \mathsf{else}\ 42/x$ remark,
        and ask which sub-expressions deserve a search rule.
\end{itemize}

\paragraph{Your turn: Task 1.} Apply the recipe to $\mathsf{true}$,
$\mathsf{false}$, and $\mathsf{if}$. Target count from the assignment: 2 value
rules, 3 step rules, exactly 1 of which is a search rule.

% --------------------------------------------------------------------
\section*{A2. Evaluating step by step \hfill (technique for Task 2)}

\paragraph{Recipe.} At each line, find the \emph{innermost-leftmost} piece
that a reduction rule can fire on. Then name the stack of search rules that
``reaches in'' to it, from the outside in.

\paragraph{Worked.} Evaluate $\FST{\PAIR{\num1+\num2}{\str{a}\,\hat{}\,\str{b}}}$:
\[
\begin{array}{r l l}
  & \FST{\PAIR{\num1+\num2}{\str{a}\,\hat{}\,\str{b}}} & \\[2pt]
  \mapsto & \FST{\PAIR{\num3}{\str{a}\,\hat{}\,\str{b}}}
          & \text{StepSearchFst $\leftarrow$ StepSearchPairL $\leftarrow$ StepAdd}\\[2pt]
  \mapsto & \FST{\PAIR{\num3}{\str{ab}}}
          & \text{StepSearchFst $\leftarrow$ StepSearchPairR $\leftarrow$ StepCat}\\[2pt]
  \mapsto & \num3
          & \text{StepFst}
\end{array}
\]
$\num3$ is a value by ValNum, so we stop. The pair $\PAIR{\num3}{\str{ab}}$
wasn't the end, because it sits inside $\FST{}$.

The ``$\leftarrow$'' stack is shorthand for a real derivation tree. Here is
the first step written out in full. Put one like this in your write-up if
you want to be extra rigorous:
\[
\Rule{StepSearchFst}
  {\Rule{StepSearchPairL}
     {\Rule{StepAdd}{\strut}{\stepto{\num1+\num2}{\num3}}}
     {\stepto{\PAIR{\num1+\num2}{\str{a}\,\hat{}\,\str{b}}}{\PAIR{\num3}{\str{a}\,\hat{}\,\str{b}}}}}
  {\stepto{\FST{\PAIR{\num1+\num2}{\str{a}\,\hat{}\,\str{b}}}}{\FST{\PAIR{\num3}{\str{a}\,\hat{}\,\str{b}}}}}
\]
The second step also needs a leaf $\val{\num3}$ (ValNum), because
StepSearchPairR has that premise.

\paragraph{Your turn: Task 2.} Your expression has three nesting levels:
$\mathsf{if}$ $\supset$ $=$ $\supset$ $+$. Your first step's stack has one
rule per level.

% --------------------------------------------------------------------
\section*{A3. Typing derivations \hfill (technique for Task 3)}

Typing rules for pairs:
\[
  \Rule{TypePair}{e_1:\tau_1 \\ e_2:\tau_2}{\PAIR{e_1}{e_2}:\tau_1\times\tau_2}
  \qquad
  \Rule{TypeFst}{e:\tau_1\times\tau_2}{\FST{e}:\tau_1}
\]

\paragraph{Recipe.} Write the goal at the bottom. Look at the \emph{outermost}
syntax to pick the rule, write its premises above the line, and repeat until
every branch ends in an axiom (a rule with no premises).

\paragraph{Worked.} $\FST{\PAIR{\num1+\num2}{\str{hi}}}:\INT$
{\small
\[
\Rule{TFst}
 {\Rule{TPair}
   {\Rule{TA}
      {\Rule{TN}{\strut}{\num1:\INT} \\ \Rule{TN}{\strut}{\num2:\INT}}
      {\num1+\num2:\INT}
    \\
    \Rule{TS}{\strut}{\str{hi}:\STRING}}
   {\PAIR{\num1+\num2}{\str{hi}}:\INT\times\STRING}}
 {\FST{\PAIR{\num1+\num2}{\str{hi}}}:\INT}
\]
}
When a rule like TypeFst lets you choose (here, any $\tau_2$), the
sub-derivation decides it. With $\mathsf{if}$, TypeIf forces both branches
to have the same $\tau$, and $\tau$ is also the type of the whole
expression.

\paragraph{Your turn: Task 3.} The outermost rule is TypeIf with
\emph{three} premises. The condition's sub-tree looks just like the
$\num1+\num2$ sub-tree above, with one more rule (TypeEq) underneath.

% --------------------------------------------------------------------
\section*{A4. Preservation cases \hfill (technique for Task 4)}

\textbf{Structure:} induction on the derivation of $\stepto{e}{e'}$, one case
per step rule. Each case follows this template:
\begin{enumerate}
  \item ``Then $e=\ldots$ and $e'=\ldots$'' (read off the rule).
  \item Invert the typing rule for $e$'s outermost form, to learn the types
        of its pieces.
  \item \emph{Reduction rule:} $e'$ is one of those pieces (or a new
        literal), so its type is already known, or comes from an axiom.\\
        \emph{Search rule:} apply the IH to the premise $\stepto{e_i}{e_i'}$,
        then re-apply the same typing rule.
\end{enumerate}

\paragraph{Worked: Case StepFst.} Then $e=\FST{\PAIR{e_1}{e_2}}$ and $e'=e_1$,
with $\val{e_1}$ and $\val{e_2}$. We know $e:\tau$. The only rule concluding
$\FST{\cdot}:\tau$ is TypeFst, so by inversion
$\PAIR{e_1}{e_2}:\tau\times\tau_2$ for some $\tau_2$. The only rule
concluding a pair's type is TypePair, so inverting again gives $e_1:\tau$ and
$e_2:\tau_2$. So $e' = e_1 : \tau$. \qed

\paragraph{Worked: Case StepSearchFst.} Then $e=\FST{e_0}$ and $e'=\FST{e_0'}$
with $\stepto{e_0}{e_0'}$. By inversion on TypeFst, $e_0:\tau\times\tau_2$ for
some $\tau_2$. By the induction hypothesis on $\stepto{e_0}{e_0'}$,
$e_0':\tau\times\tau_2$. By TypeFst, $e'=\FST{e_0'}:\tau$. \qed

\paragraph{Watch out.} In the search case, the IH is about the
\emph{sub}-expression's type ($\tau\times\tau_2$ here), not $\tau$. Say
which type you're applying it at.

\paragraph{Your turn: Task 4.} Three cases, one per Task 1 step rule. Your two
reduction cases each need \emph{one} inversion (on TypeIf).

% --------------------------------------------------------------------
\section*{A5. Canonical Forms \hfill (technique for Task 5)}

\paragraph{Recipe.} Canonical forms answers: ``which value rules could
possibly have produced a value of this type?'' List those shapes.

\paragraph{Worked.} If $\val{e}$ and $e:\tau_1\times\tau_2$, then
$e=\PAIR{e_1}{e_2}$ for some $e_1,e_2$ with $\val{e_1}$ and $\val{e_2}$.
(ValNum and ValString produce things of type $\INT$ and $\STRING$, so only
ValPair can produce a value of a product type.)

\paragraph{Your turn: Task 5.} Which of your Task 1 value rules produce
something of type $\mathsf{bool}$? A type may have more than one canonical
form, and that's fine. It just means a case split later.

% --------------------------------------------------------------------
\section*{A6. Progress cases \hfill (technique for Task 6)}

\textbf{Structure:} induction on the derivation of $e:\tau$. Template for a
rule with sub-expressions:
\begin{enumerate}
  \item Invert to get the sub-expression typings (the premises of the rule).
  \item Apply the IH to the sub-expression that gets evaluated first. You
        get two cases: it's a value, or it steps.
  \item \emph{Steps:} use the search rule.\quad
        \emph{Value:} use Canonical Forms to get its exact shape, then use
        the reduction rule for that shape.
\end{enumerate}

\paragraph{Worked: Case TypeFst.} Then $e=\FST{e_0}$ with
$e_0:\tau\times\tau_2$. By the IH, $\val{e_0}$ or $\stepto{e_0}{e_0'}$.
\begin{itemize}
  \item Suppose $\val{e_0}$. By Canonical Forms (product case),
        $e_0=\PAIR{e_1}{e_2}$ with $\val{e_1}$ and $\val{e_2}$. By StepFst,
        $\stepto{e}{e_1}$.
  \item Suppose $\stepto{e_0}{e_0'}$. By StepSearchFst,
        $\stepto{e}{\FST{e_0'}}$.
\end{itemize}
\qed

\paragraph{Worked: Case TypePair} (for the value-rule side). Here \emph{both}
components get evaluated, so we need the IH on both, nested exactly like the
class's TypeAdd case. If both are values, $\val{e}$ by ValPair (not a step!).
If $e_1$ is a value and $e_2$ steps, use StepSearchPairR. If $e_1$ steps, use
StepSearchPairL.

\paragraph{Your turn: Task 6.} TypeTrue and TypeFalse are one-liners (like
TypeNum). For TypeIf, how many sub-expressions does your semantics evaluate
first? That tells you how many IHs you need. In the value case, Canonical Forms
gives you \emph{two} shapes, so you get two sub-cases.

% ====================================================================
\newpage
\part*{Part B: Division by zero (Tasks 7--9)}

\textbf{Worked system: $\mathsf{SQ}$}, a language that prevents
\emph{square roots of negatives}. Types $\tau \bnfdef \NN \bnfalt \ANY$
(``definitely non-negative'' vs.\ ``any integer''). Expressions:
$e \bnfdef \num n \bnfalt e+e \bnfalt e-e \bnfalt e*e \bnfalt \SQRT{e}$, with
arithmetic dynamics like $\mathsf{EZ}$'s, plus
\[
  \Rule{StepSqrt}{n\ge 0}{\stepto{\SQRT{\num n}}{\num{\lfloor\sqrt n\rfloor}}}
  \qquad
  \Rule{StepSearchSqrt}{\stepto{e}{e'}}{\stepto{\SQRT{e}}{\SQRT{e'}}}
\]
So $\SQRT{\num{-4}}$ is stuck, just like $\num1/\num0$ in $\mathsf{EZ}$. As
in $\mathsf{EZ}$, if $e:\NN$ then $e:\ANY$, and Canonical Forms says
$\NN$-values are $\num n$ with $n\ge0$.

% --------------------------------------------------------------------
\section*{B1. Designing refinement rules \hfill (technique for Task 7)}

\paragraph{Recipe (the whole trick).} For each blank, try the
\emph{stronger} type first.
\begin{itemize}
  \item To show the stronger type is \textbf{unsafe}, find \emph{one}
        counterexample: inputs satisfying the premises where the output
        breaks the property.
  \item To show it's \textbf{safe}, give an arithmetic argument that works
        for \emph{all} such inputs.
\end{itemize}

\paragraph{Worked.} Base rules:
$\Rule{TypeNum}{\strut}{\num n:\ANY}$ and
$\Rule{TypeNumNN}{n\ge0}{\num n:\NN}$.
\begin{itemize}
  \item $e_1:\NN,\ e_2:\NN \Rightarrow e_1+e_2:\ ?$ \
        A sum of two non-negatives is $\ge 0$ for \emph{all} inputs. Safe, so
        $\NN$.
  \item $e_1:\NN,\ e_2:\NN \Rightarrow e_1-e_2:\ ?$ \
        Counterexample: $0-1=-1$. So only $\ANY$.
  \item $e_1:\NN,\ e_2:\NN \Rightarrow e_1*e_2:\ ?$ \
        A product of non-negatives is $\ge0$. So $\NN$.
  \item $\SQRT{e}$: StepSqrt's side condition is $n\ge0$, so the premise
        must be $e:\NN$. Result: $\lfloor\sqrt n\rfloor\ge 0$ always. So
        $\NN$.
\end{itemize}
Plus $\ANY$-versions ($e_1:\ANY,\ e_2:\ANY\Rightarrow e_1\,op\,e_2:\ANY$) so
that mixed programs still type check.

\paragraph{Your turn: Task 7.} Same drill for $\mathsf{EZ}$'s blanks. The
assignment reminds you integers can be negative, and division truncates.
Test $\num1/\num2$ before you decide the result type of TypeDiv.

% --------------------------------------------------------------------
\section*{B2. Progress with a side condition \hfill (technique for Task 8, Progress)}

\paragraph{Worked: Case TypeSqrt.} Then $e=\SQRT{e_0}$, $\tau=\NN$, and
$e_0:\NN$. By the IH, $\val{e_0}$ or $\stepto{e_0}{e_0'}$.
\begin{itemize}
  \item Suppose $\val{e_0}$. By Canonical Forms (\textbf{the $\NN$ case}),
        $e_0=\num n$ with $n\ge 0$. That's exactly StepSqrt's side condition,
        so $\stepto{e}{\num{\lfloor\sqrt n\rfloor}}$.
  \item Suppose $\stepto{e_0}{e_0'}$. By StepSearchSqrt,
        $\stepto{e}{\SQRT{e_0'}}$.
\end{itemize}
\qed

\noindent\textbf{The point:} the \emph{type} of the argument is what gives
you the side condition, through Canonical Forms. That's the whole reason
the type system exists. TypeDiv has two sub-expressions, so you'll nest two
IHs (like the class's TypeAdd case) and use the $\NZ$ case of Canonical
Forms only on the one where it matters.

% --------------------------------------------------------------------
\section*{B3. Preservation when inversion gives two rules \hfill (Task 8 \& Bonus)}

\paragraph{Worked: Case StepAdd in $\mathsf{SQ}$.} Then $e=\num{n_1}+\num{n_2}$
and $e'=\num{n_1+n_2}$. \emph{Two} typing rules conclude a sum, so inversion
gives two sub-cases:
\begin{itemize}
  \item \emph{Via TypeAddNN.} Then $\tau=\NN$ and $\num{n_1}:\NN$,
        $\num{n_2}:\NN$. Inverting again (the only rule typing a numeral
        $\NN$ is TypeNumNN): $n_1\ge0$ and $n_2\ge0$. So $n_1+n_2\ge0$, and by
        TypeNumNN, $e':\NN$.
  \item \emph{Via TypeAddAny.} Then $\tau=\ANY$. By TypeNum, $e':\ANY$.
\end{itemize}
\qed

\paragraph{Your turn.} In $\mathsf{EZ}$, only \emph{one} rule concludes a
sum, a difference, or a quotient, so StepAdd, StepSub, and StepDiv each have one
case. Ask yourself what $\tau$ inversion forces, and which base typing rule
retypes the resulting numeral. The Bonus (StepMul) is the two-rule situation
shown above.

% --------------------------------------------------------------------
\section*{B4. Finding the type system's blind spot \hfill (technique for Task 9)}

\paragraph{Recipe.} Find a rule that \emph{throws away information} (its
conclusion is weaker than what's actually true for some inputs). Build an
expression whose runtime value is fine, but only because of the information
that got thrown away. Then propose extra information to track, and explain
why the fix still blocks the bad programs.

\paragraph{Worked.} In $\mathsf{SQ}$, consider
$\SQRT{(\num{-2}*\num{-3})}$. At runtime, $(-2)(-3)=6\ge0$, so it steps to
$\num2$ with no error. But $\num{-2}$ and $\num{-3}$ can only be typed
$\ANY$, so the product is $\ANY$ and TypeSqrt rejects it. The information
``both factors are negative'' was never recorded.

\emph{Fix:} add a type $\NP$ (``non-positive'') with
\[
  \Rule{TypeNumNP}{n\le 0}{\num n:\NP}
  \qquad
  \Rule{TypeMulNP}{e_1:\NP \\ e_2:\NP}{e_1*e_2:\NN}
\]
\emph{Still safe:} each new rule is backed by a true arithmetic fact (a
product of two non-positives is $\ge 0$), so nothing that evaluates to a
negative can be typed $\NN$. $\SQRT{\num{-4}}$ is still rejected.

\paragraph{Your turn: Task 9.} Look at which $\mathsf{EZ}$ rules only ever
conclude $\MZ$. The assignment says your fix doesn't need to be complete or
practical. One or two sentences of justification is fine.

% ====================================================================
\section*{LaTeX quick reference (this homework's \texttt{macros.tex})}
\begin{itemize}
  \item \verb|\Rule{Name}{prem1 \\ prem2}{conclusion}|: a labeled rule. Use
        \verb|{\strut}| for no premises.
  \item Nest a \verb|\Rule| inside a premise to build derivation trees. Wrap
        big ones in \verb|{\small ...}| or \verb|\scriptsize|.
  \item Evaluation traces: an \verb|array| with columns
        \verb|{r l l}| (\verb|\mapsto|, expression, rule names).
  \item Overlined numerals: \verb|\overline{3}|.
\end{itemize}

\section*{Final checklist}
\begin{itemize}
  \item Task 1: 5 rules, exactly one search rule, no rule evaluates an
        un-chosen branch.
  \item Task 2: every step names its rule stack. Stop at a value.
  \item Task 3: every leaf is an axiom. Both branches get the same type.
  \item Tasks 4/6/8: say what you're inducting on. Each case names the
        rule it inverts on and the rule it concludes with.
  \item Task 8 Progress (TypeDiv): the $n_2\neq 0$ side condition is
        justified by Canonical Forms, not assumed.
  \item Task 9: an expression, why it fails to type check, a fix, and why
        the fix is still safe.
  \item Task 10: don't forget it!
\end{itemize}

\end{document}
